% Section 8: Conclusion (~700 words)

\subsection{Summary of Findings}

This paper has provided the first systematic uncertainty quantification
for algorithmic congressional redistricting. We characterized three
sources of uncertainty and derived 95\% confidence intervals for the
major performance metrics of METIS recursive bisection redistricting.

The key quantitative findings are:

\textbf{METIS seed variance} contributes CV $< 2\%$ to Polsby-Popper
for 96\% of states (48/50). The bootstrap 95\% CI for the national
mean PP is $[0.351, 0.370]$ (point estimate $0.361$, width $0.019$).
Partisan seat-share uncertainty from seeds is at most $\pm 1$ seat
for the two highest-variance states (Georgia, North Carolina), an
upper bound that is small relative to interstate variation.

\textbf{Census sampling uncertainty} contributes $\pm 0.002$ to
district population deviation at 95\% confidence, via the delta
method propagating PL 94-171 coverage error estimates. This does
not affect the validity of algorithmic population balance---it
quantifies the uncertainty of the census itself. Differential
undercount means minority VAP fractions are underestimated by
$\sim 1.6\%$ (Black) and $\sim 2.6\%$ (Hispanic), making the
majority-minority district count of 137 a \emph{conservative}
lower bound on the true count.

\textbf{Shapefile resolution measurement error} contributes
$\Delta PP < 0.005$ (mean $0.003$) across the 130$\times$ resolution
range tested in C.1. The bias direction is conservative: tract-level
PP slightly underestimates block-level PP.

\textbf{The $+22\%$ compactness improvement} is the headline finding
that algorithmic plans outperform enacted plans. The 95\% CI for this
improvement, accounting for all three uncertainty sources, is
$[+15\%, +29\%]$ for the relative improvement and $[0.044, 0.088]$
for the absolute $\Delta PP$. The improvement is positive under all
scenarios, including the worst-case lower bound.

\textbf{Efficiency gap CIs} for all 15 competitive states confirm
that the direction of algorithmic partisan lean (modest Democratic
advantage, mean $-3.0\%$) is robust to seed variance and electoral
uncertainty. The joint 95\% CI for the 15-state mean EG is
$[-4.9\%, -1.1\%]$, entirely negative. The enacted vs.\ algorithmic
gap ($\sim 8.3$ pp) is far larger than any CI width.

\subsection{Practical Recommendations}

Based on this UQ analysis, we make four practical recommendations
for researchers and expert witnesses:

\begin{enumerate}
  \item \textbf{Always report CIs.} Point estimates of PP, EG, and
    minority VAP without CIs are incomplete. Use the CIs from
    Table~\ref{tab:pp-seed-ci}, Table~\ref{tab:eg-seed-ci}, and
    the synthesis table (Table~\ref{tab:robustness}).

  \item \textbf{Distinguish sources.} Seed CI and electoral CI have
    different legal interpretations. Seed CI reflects algorithm
    implementation choice (controllable); electoral CI reflects
    turnout variation (uncontrollable). Courts may weight these
    differently.

  \item \textbf{Conservative census adjustments.} Report majority-minority
    district counts from official data, noting that true counts are
    likely higher due to differential undercount. Do not adjust
    upward without citing the specific PES estimates used.

  \item \textbf{State-specific CIs for high-variance states.} Use the
    state-specific CIs from Table~\ref{tab:pp-seed-ci} rather than
    the national mean CI for testimony about Georgia, North Carolina,
    or Wisconsin, which have higher seed variance than other states.
\end{enumerate}

\subsection{Limitations}

Two limitations of this UQ analysis should be acknowledged.

First, the bootstrap CIs from B.7 assume that 10,000 seeds constitute
an adequate sample of the seed distribution. For 43/50 non-trivial
states, convergence occurs before seed 600 (B.7, B.16), so the 10,000
sample is highly sufficient. For the three high-variance states
(GA, NC, WI) where the last improvement seed exceeds 600, the tails
of the distribution are less well-sampled.

Second, the electoral uncertainty estimation uses only three elections
(2016, 2018, 2020), all from a single redistricting cycle. True
electoral uncertainty across a full decade (ten elections: 2022, 2024,
2026, 2028, 2030) would be somewhat larger. The 2030 Census will
provide the opportunity to re-estimate these quantities with a more
complete electoral record.

\subsection{Concluding Remarks}

Uncertainty quantification is not a concession of weakness; it is a
demonstration of rigor. An expert witness who knows the 95\% CI for
their headline claim and can explain how it was computed is more
credible---not less---than one who offers only a point estimate.
The analysis in this paper provides the methodological foundation
for reporting algorithmic redistricting results with the epistemic
honesty appropriate for legal proceedings and public policy.
